\documentclass[12pt,a4paper]{article}
\usepackage[utf8]{inputenc}
\usepackage[T1]{fontenc}
\usepackage[spanish]{babel}
\usepackage[margin=1.4cm]{geometry}
\usepackage{array,tabularx}
\usepackage{enumitem}
\usepackage{titlesec}
\usepackage{xcolor}
\usepackage{hyperref}

% Restrained, high-contrast palette shared by headings, emphasis and links.
\definecolor{HeadingColor}{HTML}{243B53}
\definecolor{InstructionColor}{HTML}{704214}
\definecolor{ConceptColor}{HTML}{205C50}
\definecolor{LinkColor}{HTML}{0047AB}
\hypersetup{colorlinks=true,allcolors=LinkColor,pdfborder={0 0 0}}
\newcommand{\guidelink}[2]{\href{#1}{\textcolor{LinkColor}{\textbf{\underline{#2}}}}}
\newcommand{\instruction}[1]{\textcolor{InstructionColor}{\textbf{#1}}}
\newcommand{\concept}[1]{\textcolor{ConceptColor}{\textbf{#1}}}

\setlength{\parindent}{0pt}
\setlength{\parskip}{3pt}
\setlength{\tabcolsep}{5pt}
\renewcommand{\arraystretch}{1.12}
\setlist{leftmargin=0.6cm,topsep=3pt,partopsep=0pt,parsep=0pt,itemsep=1pt}
\titleformat{\section}{\normalsize\bfseries\color{HeadingColor}}{}{0pt}{}
\titlespacing*{\section}{0pt}{8pt}{3pt}
\newcolumntype{Y}{>{\raggedright\arraybackslash}X}
\newcommand{\checkbox}{\fbox{\rule{0pt}{0.8ex}\rule{0.8ex}{0pt}}}

\begin{document}

\begin{center}
{\Large\bfseries\color{HeadingColor} C7--C8: Definir el desafío y conectar conocimientos}\\[3pt]
Emprendimiento e Innovación --- Grado 10
\end{center}

\section*{Objetivo y punto de partida}
\concept{C7:} Definir un desafío de innovación claro, concreto y comprensible para el grupo.\\
\concept{C8:} Relacionarlo con conocimientos pertinentes de otras disciplinas.

Ya tienen el contexto y los hallazgos H de C5, y una oportunidad priorizada en C6. Ahora transformarán esa decisión en una pregunta que oriente la exploración y determinarán qué conocimientos ayudan a comprenderla. Si la selección es provisional, mantengan visible su limitación.

\section*{Documento acumulativo e integridad de la evidencia}
Esta guía es de consulta. Añadan únicamente las secciones C7 y C8 al \textbf{mismo documento del proyecto en LaTeX}; mantengan las anteriores y actualicen la fecha. Pueden continuar en \guidelink{https://www.overleaf.com/}{Overleaf}. Referencien las secciones y códigos existentes; no copien perfiles, mapas, hallazgos ni matrices. Entreguen el PDF completo actualizado.

Mantengan la distinción entre evidencia, conclusión, supuesto e idea establecida en C5--C6. La IA solo puede apoyar código, formato y organización de contenido propio; \instruction{no debe fabricar evidencia, conclusiones, decisiones, ideas ni contribuciones.} Los ejemplos son ficticios. Conserven la identificación de cualquier caso simulado autorizado.

\section*{C7 --- Transformen la oportunidad en un desafío}
Un \concept{desafío de innovación} expresa qué necesita mejorar, para quién y en qué contexto, dejando abiertas distintas respuestas. Lean la decisión de C6 y los hallazgos que la respaldan. Redacten \textbf{una pregunta}; esta lista orienta la escritura, no requiere otra tabla de respuestas:
\begin{itemize}
\item \textbf{Usuarios:} grupo específico que experimenta la necesidad.
\item \textbf{Mejora deseada:} qué necesita lograr, sin atribuir causas no demostradas.
\item \textbf{Contexto y alcance:} dónde y cuándo; qué parte de la situación abordarán.
\item \textbf{Sustento:} código de oportunidad y hallazgos H pertinentes.
\end{itemize}

\begin{center}
\concept{¿Cómo podríamos [mejora deseada] para [usuarios específicos]\\
en [contexto delimitado]?}
\end{center}

\textbf{Ejemplos ficticios:}
\begin{itemize}
\item \textbf{Muy amplio:} ¿Cómo podríamos mejorar el colegio? No delimita usuario ni necesidad.
\item \textbf{Solución anticipada:} ¿Cómo podríamos crear una aplicación de pedidos? Impone un medio.
\item \textbf{Delimitado:} ¿Cómo podríamos facilitar que los estudiantes de décimo que compran en la cafetería dispongan de tiempo para comer durante el descanso? Solo es válido si sus registros respaldan esa necesidad.
\end{itemize}
\instruction{No incluyan productos, tecnologías, causas supuestas ni metas numéricas sin respaldo.}

\section*{C7 --- Revisen y acuerden D1}
Lean la pregunta en voz alta. Cada integrante explique quién necesita qué y en qué contexto. Comprueben si sus interpretaciones coinciden, si el alcance es abordable según C6 y si la pregunta admite respuestas diversas.

Registren \textbf{solo la formulación acordada}, identificada como D1, con referencia a la oportunidad y los hallazgos H. Añadan una nota breve sobre el ajuste concreto que surgió de la revisión y su razón; si no hubo cambios, expliquen por qué. Indiquen qué queda fuera del alcance o qué sigue siendo provisional. No copien versiones completas ni vuelvan a justificar la priorización.

\concept{Evidencia de C7:} D1, su delimitación y una nota de revisión grupal. Etiqueten esta sección, por ejemplo \texttt{\textbackslash label\{sec:desafio\}}, para enlazarla después mediante \texttt{\textbackslash hyperref[sec:desafio]\{D1\}}.

\newpage
\section*{C8 --- Conecten D1 con conocimientos pertinentes}
El desafío ya está definido. Seleccionen \textbf{dos disciplinas} que aporten maneras distintas de comprenderlo: diseño, comunicación, economía, tecnología, ciudadanía u otra área pertinente. No basta con nombrarlas: identifiquen un concepto concreto y consulten un material de clase o recurso confiable.

Usen \textbf{una tabla o un diagrama} con la información siguiente. Pueden elaborarlo en \guidelink{https://app.diagrams.net/}{diagrams.net} e incorporarlo al documento; no necesitan ambos formatos.

{\small
\begin{tabularx}{\linewidth}{|p{0.08\linewidth}|Y|Y|Y|}
\hline
\textbf{Código} & \textbf{Disciplina, concepto y fuente} & \textbf{Conexión con D1} & \textbf{Aporte y límite} \\
\hline
K1 & Concepto concreto; título del material y enlace si existe. & Parte del desafío y hallazgo H que hacen pertinente el concepto. & Qué ayuda a comprender y qué permite preguntar; qué no pueden concluir todavía. \\
\hline
K2 & Otra disciplina y su fuente. & Relación específica con D1 y la evidencia. & Aporte diferente y aspecto por aprender o verificar. \\
\hline
\end{tabularx}}

\textbf{Ejemplo ficticio:} si los usuarios expresan que no comprenden la información disponible antes de comprar, la comunicación aporta criterios de claridad de mensajes para analizar esa dificultad. No demuestra que una señal sea la solución ni obliga a diseñarla.

Comprueben que cada conexión responde a un hallazgo, no a una solución que ya desean construir. Distingan el conocimiento consultado de las hipótesis sobre su contexto. Ajusten las conexiones genéricas en el mismo registro.

Cierren con \textbf{una frase de complementariedad}: expliquen qué permite comprender la combinación de K1 y K2 que una sola perspectiva dejaría incompleto. No escriban otra justificación de cada fila ni elaboren un mapa nuevo del usuario.

\concept{Evidencia de C8:} dos conexiones disciplinares sustentadas, con aportes distintos, fuentes y límites. \instruction{No desarrollen productos, presupuestos ni prototipos.}

\section*{Resultado y continuidad}
C7 aporta D1; C8 aporta K1 y K2. En C9 usarán estos códigos junto con los hallazgos H para explorar ideas diversas. En C10 examinarán cómo lo que saben del usuario orienta y limita esas ideas. Mantengan las referencias internas en el mismo documento.

\section*{Verificación final}
{\small\begin{itemize}[label=\checkbox]
\item D1 deriva de C6, precisa usuarios, necesidad y contexto, y admite distintas respuestas.
\item La nota de revisión muestra cómo el grupo aclaró la pregunta y delimitó su alcance.
\item K1 y K2 aplican conceptos concretos a D1 con hallazgos, fuentes, aportes y límites.
\item Explicamos cómo se complementan sin repetir la tabla ni anticipar una solución.
\item Conservamos C1--C6 y la integridad de la evidencia; compilamos y revisamos el PDF actualizado.
\end{itemize}}

\end{document}
